\section{The full quadratic restriction image and its residual kernel}
\label{sec:complete-quadratic-restriction}

The two cyclic-cover constructions lift every basis character of the
arithmetic quadratic target. This determines the image of the original
restriction map, while isolating its still undetermined kernel.
Write \(\nu_p=\xi_{2,p}|_N\in X_2\) for the coefficient-two
cyclic-cover lift of \(\kappa_{2,p}\), for each odd prime \(p\),
given by Theorem~\ref{thm:coefficient-two-cross-cups}.

\begin{theorem}\label{thm:complete-quadratic-restriction}
Restriction \(X_2\to Y_2\) is surjective. The exact sequence
\[
 0\longrightarrow\operatorname{Hom}_{\rm cts}(B,\F_2)
 \longrightarrow X_2\xrightarrow{\ \mathrm{res}\ }Y_2
 \longrightarrow0
\]
has an explicit \(\F_2\)-linear section \(s\), given on the
arithmetic basis by
\[
 \begin{aligned}
 s(\eta_3)&=\eta_3,& s(\psi_p)&=\psi_p,\\
 s(\kappa_{2,3})&=\tau|_N+\psi_3,&
 s(\kappa_{2,p})&=\nu_p\quad(p\ne3),\\
 s(\kappa_{\ell,q})&=\xi_{\ell,q}\quad(\ell<q\text{ odd}).
 \end{aligned}
\]
Thus, putting \(\mathcal S=s(Y_2)\),
\[
 X_2=\operatorname{Hom}_{\rm cts}(B,\F_2)\oplus\mathcal S,
 \qquad \mathcal S\xrightarrow{\ \mathrm{res}\ }Y_2
                 \text{ is an isomorphism}.
\]
This does not determine \(\operatorname{Hom}_{\rm cts}(B,\F_2)\).
\end{theorem}
\begin{proof}
The complete basis theorem for \(Y_2\) lists exactly the generators
in the displayed prescription. The arithmetic identities for
\(\eta_3\), \(\psi_p\), \(\tau\), the odd-pair lifts and the
coefficient-two cyclic-cover lifts show that every prescribed image
restricts to its basis element. Every element of \(Y_2\) is a finite
linear combination, so the prescription defines a linear section.
Surjectivity follows. The previously proved identification of the
restriction kernel with \(\operatorname{Hom}_{\rm cts}(B,\F_2)\)
gives exactness and the direct sum. Explicitly, for \(\chi\in X_2\),
\(\chi-s(\mathrm{res}\chi)\) is the unique component in that kernel.
No assertion about its vanishing is used.
\end{proof}

\begin{corollary}\label{cor:full-quadratic-inflation-kernel}
The two quadratic inflation kernels on the arithmetic quotient agree:
\[
 \ker\bigl(H^2(A,\F_2)\to H^2(G,\F_2)\bigr)
 =\ker\bigl(H^2(A,\F_2)\to\operatorname{Br}(\Q)[2]\bigr).
\]
\end{corollary}
\begin{proof}
The two five-term exact sequences identify these kernels with the
transgression images of \(X_2\) and \(Y_2\), respectively.
Naturality identifies the arithmetic transgression of a restriction
with the same class in \(H^2(A,\F_2)\). Surjectivity of restriction
therefore identifies the two images.
\end{proof}

The same lifts give a further simultaneous arithmetic normalization.
Let \(\mathcal J\) be the set of odd primes other than three, and put
\[
 \begin{split}
 R_{\rm quad}&=R^*\cap\bigcap_{p\in\mathcal J}\ker\nu_p,\\
 V_{\rm quad}&=V^*\cap R_{\rm quad},\qquad
 \mathcal Q=\prod_{p\in\mathcal J}\F_2.
 \end{split}
\]
Corollary~\ref{cor:cubic-two-lifts-equal} proves
\(\nu_3=\tau+\psi_3\) on \(N\). Writing the lift in this
form displays its relation to the entire cubic logarithmic coordinate.

\begin{theorem}\label{thm:complete-quadratic-normalization}
The joint map \((\nu_p)_{p\in\mathcal J}:R^*\to\mathcal Q\)
is onto already on \(R^*\cap H\), and also on \(V^*\cap H\).
Its full and arithmetic kernels are the corresponding intersections
with \(R_{\rm quad}\). In particular
\[
 R^*/R_{\rm quad}\simeq V^*/V_{\rm quad}\simeq\mathcal Q,
 \qquad G=R_{\rm quad}H.
\]
Moreover
\[
 \begin{aligned}
 B&\simeq R_{\rm quad}/
            \overline{[G,R_{\rm quad}](R_{\rm quad}\cap H)},\\
 B_{\max2}&\simeq
 \left(V_{\rm quad}/
            \overline{[G,V_{\rm quad}](V_{\rm quad}\cap H)}\right)_{\max2}.
 \end{aligned}
\]
The full joint quotients, also on arithmetic intersections, are
\[
 R_D/R_{\rm quad}\simeq V_D/V_{\rm quad}
       \simeq\mathcal U\times\mathcal P\times\mathcal Q.
\]
The kernel of restriction from \(X_2\) to either residual subgroup
is exactly \(\mathcal S=s(Y_2)\). Its image is exactly the space
of invariant quadratic characters on that subgroup vanishing on
its arithmetic intersection, canonically identified with
\(\operatorname{Hom}_{\rm cts}(B,\F_2)\).
\end{theorem}
\begin{proof}
The preceding simultaneous theorem gives \(G=R^*H\), with original
restriction kernel
\[
 \Sigma_0\oplus\F_2\tau\oplus
                       \bigoplus_{\mathcal I}\F_2\xi_{\ell,q}.
\]
By Lemma~\ref{lem:arithmetic-normalization-annihilator}, its arithmetic
image is exactly the kernel of restriction from \(Y_2\) to
\(R^*\cap H\). The complete basis of \(Y_2\) therefore makes the
remaining \(\kappa_{2,p}\), \(p\in\mathcal J\), independent on
this arithmetic subgroup. Finite linear character separation and
compactness, exactly as in the preceding theorem, prove
\[
               (\nu_p)_{\mathcal J}(R^*\cap H)=\mathcal Q.
\]
The quotient \((R^*\cap H)/(V^*\cap H)\) embeds in the pro-odd
group \(R^*/V^*\). Hence its quotient in the pro-two image
\(\mathcal Q\) is trivial, proving the full arithmetic image on
\(V^*\cap H\). Matching these values proves
\(G=R_{\rm quad}H\), the two first quotient descriptions, and
preservation of the full relative abelian quotient.

To prescribe all three factors of
\(\mathcal U\times\mathcal P\times\mathcal Q\), first choose an
arithmetic element with the specified \(\mathcal U\times\mathcal P\)
coordinates using Theorem~\ref{thm:simultaneous-quadratic-normalization}.
An element of \(R^*\cap H\), or of \(V^*\cap H\), corrects its
\(\mathcal Q\)-coordinate without changing the first two. This proves
the full joint quotient assertion and its arithmetic version.
The quotient \(R_{\rm quad}/V_{\rm quad}\) embeds in the pro-odd
quotient \(R_D/V_D\). Applying the relative pro-odd theorem with
\(G=R_{\rm quad}H\) gives the stated maximal pro-two quotient.

An original character vanishing on \(R_{\rm quad}\) restricts to
a continuous character of the product
\(R_D/R_{\rm quad}\). Its quadratic character space is the direct
sum of the \(\tau\)-line and finitely supported coordinates of
\(\mathcal P\) and \(\mathcal Q\). Subtracting these lifts leaves
exactly \(\Sigma_0\). This is precisely \(\mathcal S\), since
replacing \(\tau\) by \(\tau+\psi_3\) preserves the span.
The pro-odd quotient gives the same kernel on \(V_{\rm quad}\).

For any \(\chi\in X_2\), subtract \(s(\mathrm{res}\chi)\).
The subtracted character vanishes on both residual subgroups, and
the difference vanishes on \(K\). Thus every restricted character
is invariant and zero on the arithmetic intersection. Conversely
such a character on \(R_{\rm quad}\) extends by zero on \(H\)
using \(G=R_{\rm quad}H\). On \(V_{\rm quad}\), the relative
pro-odd theorem first gives that same extension. Restriction of the
resulting character of \(G\) to \(N\) lies in \(X_2\).
This proves the exact image assertion and its canonical
identification with the relative kernel.
\end{proof}
